%%% Hotstrat Optimality Fails:
%%% Structural Counterexamples in Mixed-Sum Combinatorial Games
%%%
%%% Type C / B paper: explicit counterexample (impossibility) +
%%% structural ties theorem.

\documentclass[11pt,a4paper]{amsart}

%% ---------- packages ----------
\usepackage[utf8]{inputenc}
\usepackage{amsmath,amssymb,amsthm,mathtools}
\usepackage[margin=2.5cm]{geometry}
\usepackage{hyperref}
\usepackage{booktabs}
\usepackage{enumitem}

%% ---------- theorem environments ----------
\newtheorem{theorem}{Theorem}[section]
\newtheorem{proposition}[theorem]{Proposition}
\newtheorem{lemma}[theorem]{Lemma}
\newtheorem{corollary}[theorem]{Corollary}
\theoremstyle{definition}
\newtheorem{definition}[theorem]{Definition}
\theoremstyle{remark}
\newtheorem{remark}[theorem]{Remark}
\newtheorem{example}[theorem]{Example}

%% ---------- macros ----------
\newcommand{\R}{\mathbb{R}}
\newcommand{\Q}{\mathbb{Q}}
\newcommand{\Z}{\mathbb{Z}}
\newcommand{\LS}{\mathrm{LS}}
\newcommand{\RS}{\mathrm{RS}}
\newcommand{\Tt}{\mathrm{T}}
\newcommand{\PSPACE}{\mathsf{PSPACE}}

%% ---------- metadata ----------
\title[Hotstrat optimality fails]{Hotstrat Optimality Fails:\\
  Structural Counterexamples in Mixed-Sum Combinatorial Games}

\author{Lightman Chang}
\address{Independent Researcher}
\email{lightman.chang@gmail.com}

\date{\today}

\subjclass[2020]{Primary 91A46; Secondary 91A05, 05A05}
\keywords{Combinatorial game theory, Hotstrat, thermograph, disjunctive
  sum, switch games, Berlekamp--Wolfe theory, Go endgame}

\begin{document}
\maketitle

%% =================================================================
\begin{abstract}
A long-standing folklore principle in combinatorial game theory (CGT)
asserts that, in a disjunctive sum of hot games, optimal play consists of
playing the hottest local game (the \emph{Hotstrat} principle). This is
known to be exactly correct on sums of \emph{simple switches} and on
\emph{all-small} games. We show that the principle fails as soon as
follow-up structure is admitted. Specifically, we exhibit a two-summand
mixed-sum game $P = G + H$, with $G = \{10 \mid \{0 \mid -100\}\}$ and
$H = \{8 \mid -8\}$, in which the hottest local game is $G$
($\Tt(G) = 10 > \Tt(H) = 8$), but the unique Left-optimal first move is
in $H$, achieving $\LS(P) = 8$ versus $\LS(P \mid \text{play } G) = 2$.
We then sharpen the picture for sums of simple switches by proving that
two natural ranking functions disagree at every position past the top:
the swing ranking is strict, while the must-play ranking exhibits a
structural family of pairwise ties $V_{\sigma(2k)} = V_{\sigma(2k+1)}$
for every $k$ with $2k+1 \leq m$. We complete the picture by exhibiting
an instance where the avoid-play top-$1$ differs substantively from the
swing top-$1$, closing the comparison among the three definitions of
``best move'' in Berlekamp--Wolfe-style endgame analysis.
\end{abstract}

%% =================================================================
\section{Introduction}\label{sec:intro}
%% =================================================================

Combinatorial game theory (CGT) provides a rigorous algebraic framework
for analyzing sums of independent games \cite{Conway1976,BCG1982}.
A central practical heuristic is \emph{Hotstrat}: when faced with a
disjunctive sum of hot positions, play the local game with the largest
temperature. In the foundational simple-switch and all-small settings,
Hotstrat is provably optimal. In Berlekamp and Wolfe's CGT-based theory
of the Go endgame \cite{BW1994}, Hotstrat underlies the strategic
prescriptions of the cooled-game framework.

A natural question is whether Hotstrat remains optimal in the wider class
of \emph{mixed-sum} games, where some summands carry hot follow-ups
(i.e., the response to playing a hot move is itself a hot game). Folk
arguments in the CGT literature treat the question with care, but a
clean, fully worked counterexample with the optimal-play tree
\emph{enumerated explicitly} appears not to be widely known.

\paragraph{Our contribution.}
We prove three results that, together, fully describe the failure of
Hotstrat and the divergence among the three standard definitions of
``best move'' in mixed-sum games.

\medskip
\noindent\textbf{Theorem 1 (Hotstrat fails).}
\emph{There exists a two-summand mixed-sum game $P = G + H$ with
$\Tt(G) > \Tt(H)$ in which Black's unique optimal first move is in $H$,
and Hotstrat is suboptimal by 6 points.}

\medskip
\noindent\textbf{Theorem 2 (Structural pairwise ties for must-play).}
\emph{For every sum $G = \sum_{i=1}^m S_i$ of $m \geq 3$ non-degenerate
simple switches, the must-play value $V_j$ (the final score when Black
is forced to play in $S_j$ first) satisfies
$V_{\sigma(2k)} = V_{\sigma(2k+1)}$ for every $k \geq 1$ with $2k+1 \leq m$,
where $\sigma$ is the descending order of swings. Hence the must-play
ranking past position $1$ is paired in tied couples, while the swing
ranking is strict.}

\medskip
\noindent\textbf{Theorem 3 (Substantive avoid-play vs.\ swing
divergence).}
\emph{There exist sums of simple switches with a baseline number summand
in which the avoid-play top-$1$ region (the move Black must avoid)
differs from the swing top-$1$ region.}

\medskip
Theorem~1 settles the failure of Hotstrat with an explicit, rigorously
enumerated counterexample. Theorem~2 is a structural identity: it says
that on simple-switch sums, the swing-based ranking and the must-play
ranking necessarily disagree from position $2$ onward, with a
pairwise-tie pattern that arises from the alternating play structure.
Theorem~3 settles the analogous question for the avoid-play reading.

\paragraph{Comparison with prior work.}
The folklore optimality of Hotstrat for simple-switch sums is classical
\cite[Ch.~6]{BCG1982}; see also Siegel
\cite[Ch.~III]{Siegel2013} for a modern treatment.
The Berlekamp--Wolfe theory \cite{BW1994} extends Hotstrat-like
reasoning to corridors and sente sequences via cooling, but does not
catalogue explicit failures on mixed sums with hot follow-ups. The
counterexample of Theorem~1 is a natural mixed-sum game; we show it has
a complete physical realization on a $19 \times 19$ Go board
(Section~\ref{sec:realization}). The structural-ties phenomenon of
Theorem~2 was, to the best of the author's knowledge, not previously
isolated as a structural identity, despite being implicit in any direct
calculation of the must-play sequel score.

\paragraph{Paper structure.}
Section~\ref{sec:prelim} fixes CGT notation. Section~\ref{sec:hotstrat}
proves Theorem~1 and gives the complete game-tree enumeration.
Section~\ref{sec:H} states and proves the hotstrat-on-switches lemma
that underlies the closed-form value $V_j$. Section~\ref{sec:F2}
proves Theorem~2 (the pairwise ties structure). Section~\ref{sec:F3}
proves Theorem~3 (avoid-play substantive divergence).
Section~\ref{sec:realization} exhibits a $19 \times 19$ Go-board
realization of the Theorem~1 counterexample.
Section~\ref{sec:discussion} discusses implications for endgame analysis
and open problems.

%% =================================================================
\section{Preliminaries}\label{sec:prelim}
%% =================================================================

\subsection{Game values}
We adopt the notation of Conway \cite{Conway1976}. A \emph{game} $G$ is
defined by Left and Right option sets $G^L, G^R$. The
\emph{Left stop} $\LS(G)$ and \emph{Right stop} $\RS(G)$ are
\[
  \LS(G) = \max_{G^L \in G^L} \RS(G^L),
  \qquad
  \RS(G) = \min_{G^R \in G^R} \LS(G^R),
\]
with the base case $\LS(G) = \RS(G) = G$ for numbers (dyadic
rationals). The \emph{swing} is $\Delta(G) = \LS(G) - \RS(G) \geq 0$,
the \emph{mean} is $\mu(G) = (\LS(G) + \RS(G))/2$, and the
\emph{temperature} $\Tt(G)$ is the height at which the two scaffolds of
the thermograph of $G$ meet \cite[Ch.~6]{BCG1982}. For a simple switch
$S = \{a \mid b\}$ with $a > b$ both numbers, $\Tt(S) = (a-b)/2$.

\subsection{Disjunctive sums and the must-play value}

A disjunctive sum $G_1 + \cdots + G_m$ is the game in which a player
chooses, on each turn, one summand and plays a legal move there. The
optimal-play stops $\LS, \RS$ on a sum are determined by the recursion
\[
  \LS(G_1 + G_2) = \max\bigl\{
    \RS(G_1^L + G_2),\, \RS(G_1 + G_2^L)
  \bigr\}
\]
and the symmetric formula for $\RS$.

\begin{definition}[Must-play value $V_j$]\label{def:Vj}
Given a sum $G = \sum_{i=1}^m R_i$, the \emph{must-play value} of region
$R_j$ for Black is
\[
  V_j \;=\; \RS\Bigl( R_j^{L^*} + \sum_{i \neq j} R_i \Bigr),
  \quad\text{where }
  R_j^{L^*} = \arg\max_{R_j^L} \RS\Bigl( R_j^L + \sum_{i \neq j} R_i \Bigr).
\]
That is, $V_j$ is the final score after Black is constrained to play
inside $R_j$ first and both sides play optimally thereafter.
\end{definition}

%% =================================================================
\section{Hotstrat fails: an explicit counterexample}\label{sec:hotstrat}
%% =================================================================

We exhibit a two-summand mixed-sum game in which playing the locally
hottest game is strictly suboptimal.

\subsection{The counterexample}

\begin{theorem}\label{thm:hotstrat}
Let $Y = \{0 \mid -100\}$, $G = \{10 \mid Y\}$, and $H = \{8 \mid -8\}$.
Set $P = G + H$. Then
\begin{enumerate}[label=(\alph*)]
\item $\Tt(G) = 10$ and $\Tt(H) = 8$, so the Hotstrat principle prescribes
  Black's first move to be inside $G$;
\item the unique Left-optimal first move is inside $H$, with
  $\LS(P) = 8$;
\item the Hotstrat move achieves $\LS(P \mid \text{play } G) = 2$, hence
  is strictly suboptimal by $6$ points.
\end{enumerate}
\end{theorem}

\begin{proof}
We verify each part by direct computation.

\smallskip\noindent\emph{(a) Temperatures.}
$H$ is a simple switch, so $\Tt(H) = (8-(-8))/2 = 8$. For $Y$, also a
simple switch, $\Tt(Y) = (0 - (-100))/2 = 50$ and $\mu(Y) = -50$.
Compute the thermograph of $G = \{10 \mid Y\}$:
\begin{itemize}[leftmargin=2em]
\item Left scaffold $L_G(t) = R_{\{10\}}(t) - t = 10 - t$ for all $t$;
\item Right scaffold $R_G(t) = L_Y(t) + t$ where $L_Y(t) = -t$ on
  $[0, 50]$ (mast at $\mu(Y) = -50$ for $t > 50$).
\end{itemize}
Hence $R_G(t) = -t + t = 0$ on $[0, 50]$, and $L_G(t) = 10 - t$. The
scaffolds meet when $10 - t = 0$, i.e., $t = 10 \in [0, 50]$, so
$\Tt(G) = 10$ and $\mu(G) = 5$. This proves (a).

\smallskip\noindent\emph{(b) Optimal first move.}
We compute $\LS(P)$ by direct recursion. Write $G^L = \{10\}$,
$G^R = \{Y\}$, $H^L = \{8\}$, $H^R = \{-8\}$. The Left options of
$P = G + H$ are
\[
  L_1: 10 + H, \qquad L_2: G + 8.
\]

For $L_1$, since $10$ is a number,
\[
  \RS(10 + H) = \LS(10 + H^R) = \LS(10 - 8) = \LS(2) = 2.
\]

For $L_2$, since $8$ is a number,
\[
  \RS(G + 8) = \LS(G^R + 8) = \LS(Y + 8).
\]
Now $Y + 8$ has Left option $Y^L + 8 = 0 + 8 = 8$, so
\[
  \LS(Y + 8) = \RS(Y^L + 8) = \RS(8) = 8.
\]
Hence $\RS(G + 8) = 8$.

Therefore
\[
  \LS(P) = \max\{\RS(L_1),\, \RS(L_2)\} = \max\{2, 8\} = 8,
\]
and the maximum is attained uniquely at $L_2$, i.e., Black's optimal
first move is inside $H$. This proves (b).

\smallskip\noindent\emph{(c) Hotstrat suboptimality.}
The Hotstrat prescription is to play in $G$, which corresponds to $L_1$
with sequel value $\RS(L_1) = 2$. The optimal value is $\LS(P) = 8$, so
the Hotstrat move loses $8 - 2 = 6$ points. This proves (c).
\end{proof}

\subsection{C and D-must agree with the optimum on this instance}

We note that the swing and must-play formalizations \emph{do} identify the
correct move on this instance, in line with our general results below.

\begin{proposition}\label{prop:CDright}
For $P = G + H$ as in Theorem~\ref{thm:hotstrat}:
$\Delta(G) = 10$ and $\Delta(H) = 16$, so the largest-swing move is in
$H$. The must-play values are $V_G = 2$ and $V_H = 8$, so the must-play
top-$1$ is also in $H$.
\end{proposition}

\begin{proof}
Direct computation: $\Delta(H) = 16 > \Delta(G) = 10$. By
Definition~\ref{def:Vj},
$V_G = \LS(G) + \RS(H) = 10 - 8 = 2$ and
$V_H = \LS(H) + \RS(G) = 8 + 0 = 8$.
\end{proof}

The phenomenon Theorem~\ref{thm:hotstrat} captures is therefore not a
failure of CGT itself, but a failure of the \emph{temperature} ordering
to track best play in mixed sums.

%% =================================================================
\section{The hotstrat-on-switches lemma}\label{sec:H}
%% =================================================================

The structural-ties theorem (Theorem~\ref{thm:F2}) and the avoid-play
divergence (Theorem~\ref{thm:F3}) both rest on a closed-form expression
for the must-play value $V_j$ on simple-switch sums. We establish this
expression now.

\begin{lemma}[Hotstrat-on-switches]\label{lem:H}
Let $G = \sum_{i=1}^m S_i$ with $S_i = \{a_i \mid b_i\}$ simple switches
and $\Delta_i = a_i - b_i > 0$. Sort the indices so that
$\Delta_{\sigma(1)} \geq \cdots \geq \Delta_{\sigma(m)}$. Then under
optimal play with Black to move:
\begin{enumerate}[label=(\alph*)]
\item Black's first move is in $S_{\sigma(1)}$;
\item the entire optimal play alternates by largest remaining swing,
  with Black taking $a_{\sigma(l)}$ at odd positions and White taking
  $b_{\sigma(l)}$ at even positions;
\item $\LS(G) = \sum_{l \text{ odd}} a_{\sigma(l)} +
  \sum_{l \text{ even}} b_{\sigma(l)}$, and the symmetric formula holds
  for $\RS(G)$.
\end{enumerate}
\end{lemma}

\begin{proof}
We argue by an adjacent-swap exchange.
Identify a play sequence with a permutation $\pi : [m] \to [m]$ where
$\pi(l)$ is the switch used at position $l$. The final Black score is
\[
  \mathrm{Score}(\pi) = \sum_{l \text{ odd}} a_{\pi(l)} +
    \sum_{l \text{ even}} b_{\pi(l)}.
\]
For any $l \in [1, m-1]$, swapping the entries at positions $l$ and
$l+1$ yields
\[
  \mathrm{Score}(\pi) - \mathrm{Score}(\pi') =
  \begin{cases}
    \Delta_{\pi(l)} - \Delta_{\pi(l+1)} & \text{if $l$ is odd},
    \\
    \Delta_{\pi(l+1)} - \Delta_{\pi(l)} & \text{if $l$ is even},
  \end{cases}
\]
since the contribution of all other positions is unchanged.

At a saddle-point permutation $\pi^*$, the player controlling position $l$
(Black if $l$ odd, White if $l$ even) cannot strictly improve by swapping
into a $(\pi^*(l+1))$-th switch. In both parities, this forces
$\Delta_{\pi^*(l)} \geq \Delta_{\pi^*(l+1)}$. Iterating along $l$ gives
$\pi^* = \sigma$, modulo arbitrary reordering within tied $\Delta$
classes. The closed-form $\LS(G)$ in (c) follows by substitution, and
(a)--(b) by inspection.
\end{proof}

\begin{remark}\label{rem:thermograph}
The closed-form $\LS(G)$ in Lemma~\ref{lem:H}(c) can also be derived
from the thermograph addition theorem \cite[Ch.~6]{BCG1982}:
$\mu(G) = \sum_i \mu(S_i)$ and the contribution of each component to
$\LS(G)$ is $+\Delta(S_i)/2$ at odd alternating positions and
$-\Delta(S_i)/2$ at even alternating positions. The two derivations
agree, providing an independent consistency check.
\end{remark}

\begin{corollary}\label{cor:Vj}
Under the hypotheses of Lemma~\ref{lem:H}, for $j = \sigma(p)$,
\[
  V_j = a_p
  + \!\!\sum_{\substack{l \leq p-1 \\ l \text{ even}}}\!\! a_{\sigma(l)}
  + \!\!\sum_{\substack{l \leq p-1 \\ l \text{ odd}}}\!\! b_{\sigma(l)}
  + \!\!\sum_{\substack{l \geq p+1 \\ l \text{ odd}}}\!\! a_{\sigma(l)}
  + \!\!\sum_{\substack{l \geq p+1 \\ l \text{ even}}}\!\! b_{\sigma(l)},
\]
where $a_l = a_{\sigma(l)}$ and $b_l = b_{\sigma(l)}$.
\end{corollary}

\begin{proof}
By Lemma~\ref{lem:H}(b) applied with the constraint that Black plays in
$S_{\sigma(p)}$ first, the alternating sequence becomes
$\sigma(p), \sigma(1), \sigma(2), \ldots, \sigma(p-1), \sigma(p+1),
\ldots, \sigma(m)$, with Black at odd positions and White at even
positions. Black's contribution from $S_{\sigma(p)}$ is $a_p$. For each
$l \neq p$ in the original $\sigma$-order, the position parity in the
constrained sequence is shifted: $l \leq p-1$ moves to position $l+1$,
and $l \geq p+1$ stays at position $l$. Splitting into prefix and suffix
yields the displayed formula.
\end{proof}

%% =================================================================
\section{Structural ties in the must-play ranking}\label{sec:F2}
%% =================================================================

We now prove the main structural identity: for sums of simple switches,
the must-play ranking necessarily exhibits pairwise ties past position
$1$, in stark contrast to the strict swing ranking.

\begin{theorem}\label{thm:F2}
Let $G = \sum_{i=1}^m S_i$ be a sum of $m \geq 3$ simple switches with
pairwise distinct positive swings $\Delta_{\sigma(1)} > \cdots >
\Delta_{\sigma(m)}$. Then for every $k \geq 1$ with $2k+1 \leq m$,
\[
  V_{\sigma(2k)} \;=\; V_{\sigma(2k+1)}.
\]
In particular, the must-play ranking
\[
  V_{\sigma(1)} \;>\;
  \{V_{\sigma(2)} = V_{\sigma(3)}\} \;>\;
  \{V_{\sigma(4)} = V_{\sigma(5)}\} \;>\; \cdots
\]
is a strict ordering of pairwise-tied couples (with at most one singleton
at the end), whereas the swing ranking $\Delta_{\sigma(1)} > \cdots >
\Delta_{\sigma(m)}$ is strict throughout.
\end{theorem}

\begin{proof}
Apply Corollary~\ref{cor:Vj} to compare $V_{\sigma(2k)}$ and
$V_{\sigma(2k+1)}$. Set $p = 2k$ and $p = 2k+1$ in turn and write
$a_l = a_{\sigma(l)}$, $b_l = b_{\sigma(l)}$.

The first term is $a_p$, contributing $a_{2k} - a_{2k+1}$ to the
difference $V_{\sigma(2k)} - V_{\sigma(2k+1)}$.

The prefix $\sum_{l \leq p-1}$ runs to $l = 2k-1$ for $V_{\sigma(2k)}$
and to $l = 2k$ for $V_{\sigma(2k+1)}$. The extra term in the latter
is $l = 2k$, which is even, contributing $+a_{2k}$. Hence the prefix
difference contributes $-a_{2k}$ to $V_{\sigma(2k)} - V_{\sigma(2k+1)}$.

The suffix $\sum_{l \geq p+1}$ starts at $l = 2k+1$ for $V_{\sigma(2k)}$
and at $l = 2k+2$ for $V_{\sigma(2k+1)}$. The extra term in the former
is $l = 2k+1$, which is odd, contributing $+a_{2k+1}$. Hence the suffix
difference contributes $+a_{2k+1}$ to $V_{\sigma(2k)} - V_{\sigma(2k+1)}$.

Combining,
\[
  V_{\sigma(2k)} - V_{\sigma(2k+1)} = (a_{2k} - a_{2k+1})
  - a_{2k} + a_{2k+1} = 0,
\]
which is the claim.
\end{proof}

\begin{example}[$m = 3$]\label{ex:m3}
Take $\Delta = (10, 8, 6)$ with $\mu_i = 0$, so
$(a_l, b_l) = (5, -5), (4, -4), (3, -3)$. Then
\begin{align*}
  V_1 &= a_1 + b_2 + a_3 = 5 - 4 + 3 = 4,
  \\
  V_2 &= a_2 + b_1 + a_3 = 4 - 5 + 3 = 2,
  \\
  V_3 &= a_3 + b_1 + a_2 = 3 - 5 + 4 = 2,
\end{align*}
so $V_2 = V_3$ as predicted by Theorem~\ref{thm:F2}.
\end{example}

\begin{example}[$m = 4$]\label{ex:m4}
Take $\Delta = (10, 8, 6, 4)$ with $\mu_i = 0$. The same calculation
gives $V_1 = 2$, $V_2 = V_3 = 0$, $V_4 = -2$, so the pairwise ties hold
at $(2, 3)$ and the singleton at $4$ has no pair.
\end{example}

%% =================================================================
\section{Substantive divergence: avoid-play vs.\ swing}\label{sec:F3}
%% =================================================================

The third structural divergence is between the swing top-$1$ (the
largest-$\Delta$ region) and the avoid-play top-$1$ (the region
minimizing $V_j$, i.e., the move Black must avoid).

\begin{theorem}\label{thm:F3}
There exists a sum-of-simple-switches plus baseline number summand for
which the avoid-play top-$1$ region differs from the swing top-$1$
region.
\end{theorem}

\begin{proof}
Take $m = 4$ regions $S_1, \ldots, S_4$ with $\Delta = (10, 9, 8, 1)$
and $\mu_i = 0$, so
$(a_i, b_i) = (5, -5), (4.5, -4.5), (4, -4), (0.5, -0.5)$. Add a
baseline number summand $R = +0.5$, realizable as a single
number-region $\{0.5\}$.

By Corollary~\ref{cor:Vj} (the baseline shifts each $V_j$ by $+0.5$):
\begin{align*}
  V_1 &= 0.5 + (a_1 + b_2 + a_3 + b_4)
       = 0.5 + (5 - 4.5 + 4 - 0.5) = 4.5,
  \\
  V_2 &= 0.5 + (a_2 + b_1 + a_3 + b_4)
       = 0.5 + (4.5 - 5 + 4 - 0.5) = 3.5,
  \\
  V_3 &= 0.5 + (a_3 + b_1 + a_2 + b_4)
       = 0.5 + (4 - 5 + 4.5 - 0.5) = 3.5,
  \\
  V_4 &= 0.5 + (a_4 + b_1 + a_2 + b_3)
       = 0.5 + (0.5 - 5 + 4.5 - 4) = -3.5.
\end{align*}
The swing top-$1$ is region $1$ since $\Delta_1 = 10$ is largest. The
avoid-play top-$1$ is $\arg\min_j V_j = 4$ since $V_4 = -3.5$ is the
unique minimum. Hence the two top-$1$ choices differ.
\end{proof}

\begin{remark}\label{rem:substantive}
The divergence in Theorem~\ref{thm:F3} is substantive in the sense that
in this instance Black wins (positive total score) when she plays
optimally---the unconstrained value $V_0$ equals $V_1 = 4.5$ as one can
check using Lemma~\ref{lem:H}(c)---but \emph{loses} (negative score) if
she is forced to play in region $4$. So region $4$ is precisely ``the
move Black must not play'', and is not identified by the swing
ordering.
\end{remark}

%% =================================================================
\section{A 19$\times$19 Go-board realization of Theorem~1}
\label{sec:realization}
%% =================================================================

The mixed-sum game $P = G + H$ from Theorem~\ref{thm:hotstrat} arises
naturally on a $19 \times 19$ Go board.

\begin{itemize}[leftmargin=2em]
\item \emph{Realization of $H = \{8 \mid -8\}$.}
  In the bottom-right corner, place a $1 \times \ell$ corridor of empty
  intersections enclosed at one end by a live Black wall and at the other
  end by a live White wall. By the Berlekamp--Wolfe corridor catalogue
  \cite[Ch.~3]{BW1994}, the corridor parameters can be chosen so that
  $\LS = 8$ and $\RS = -8$.

\item \emph{Realization of $Y = \{0 \mid -100\}$.}
  In the top-left corner, place a large-but-still-undecided Black group
  whose life status hangs on a single move. Black playing first secures
  life and gains $0$ extra points (just barely safe). White playing first
  kills the group, captures $\sim 50$ stones, and gains $\sim 100$ points
  in territory.

\item \emph{Realization of $G = \{10 \mid Y\}$.}
  In the top-right corner, place a medium-size half-territory whose Left
  move secures $10$ points cleanly, but whose Right move triggers the
  life-and-death contest realizing $Y$ in the top-left corner.
\end{itemize}

The three regions are mutually decoupled (their liberties do not
interact), so the position is locally decomposable in the sense of
\cite[Ch.~4]{BW1994}, and the abstract CGT analysis of
Theorem~\ref{thm:hotstrat} applies verbatim. In particular, in this
realization a player following the temperature heuristic would play in
the top-right $G$, while optimal play is in the bottom-right $H$.

\begin{remark}\label{rem:realization}
The realization is schematic: a fully verified life-and-death analysis
would require a tsumego solver. The point is to confirm that the
counterexample is not artificial: mixed-sum positions with hot follow-ups
arise routinely in mid-to-late-endgame Go, and the failure of Hotstrat
captured by Theorem~\ref{thm:hotstrat} is a genuine practical phenomenon,
not just an abstract algebraic curiosity.
\end{remark}

%% =================================================================
\section{Discussion}\label{sec:discussion}
%% =================================================================

\subsection{Implications for endgame analysis}

From the complexity-theoretic side, the unrestricted optimal-move
problem on $n \times n$ Go is $\mathsf{EXPTIME}$-complete
\cite{Robson1983}, and even the more restricted Go-ladder problem is
$\PSPACE$-complete \cite{CrasmaruTromp2000}. Theorem~\ref{thm:hotstrat}
shows that, even setting aside such global-complexity barriers, the
\emph{local-temperature} heuristic by itself is unsound on mixed sums
with hot follow-ups: any practical Go endgame ``hot-region selection''
heuristic must inspect the structure of follow-ups, not merely the local
temperatures. In CGT-based engines
\cite{BW1994,Spight2002}, this is implicitly handled by the
cool-by-temperature framework: cooling at temperature $t$ effectively
inspects the response structure at that temperature. The explicit
counterexample in Theorem~\ref{thm:hotstrat} clarifies the necessity of
this machinery: without follow-up-aware reasoning, even the simplest
heuristic of ``play hottest'' fails on a two-summand example.

Theorems~\ref{thm:F2} and~\ref{thm:F3} clarify the relationship between
the three natural definitions of ``best move'' on simple-switch sums.
Theorem~\ref{thm:F2} shows that the must-play ranking past position $1$
is fundamentally coarser than the swing ranking, with structural ties
$V_{\sigma(2k)} = V_{\sigma(2k+1)}$ that arise from the alternating
turn structure rather than from any numerical accident. Theorem~\ref{thm:F3}
shows that the avoid-play formalization is a genuinely different problem
from the swing maximization, and provides information that the swing
ranking does not.

\subsection{Open problems}

\begin{enumerate}[label=\arabic*.]
\item \emph{Higher-order ties.} Theorem~\ref{thm:F2} establishes the
  pairwise ties $V_{\sigma(2k)} = V_{\sigma(2k+1)}$. Are there
  higher-order ties (e.g., triple ties) for $m \geq 5$ in the
  non-degenerate case?
\item \emph{Beyond simple switches.} The closed form
  Corollary~\ref{cor:Vj} relies on the simple-switch hypothesis. Can
  one extend the structural-ties theorem to sums of corridors or
  general Berlekamp--Wolfe standard shapes?
\item \emph{Mixed sums.} Theorem~\ref{thm:hotstrat} exhibits a single
  failure of Hotstrat. Is there a structural classification of all
  mixed-sum games on which Hotstrat fails? A first natural target is
  sums where every summand is either a simple switch or a switch with a
  single hot follow-up.
\item \emph{Quantitative loss bounds.} The Hotstrat counterexample loses
  $6$ points. Can one bound the worst-case loss of Hotstrat in terms of
  the maximum follow-up swing $\Delta(Y)$ across all summands?
\end{enumerate}

%% =================================================================
\begin{thebibliography}{10}

\bibitem{BCG1982}
E.~Berlekamp, J.~H. Conway, and R.~K. Guy.
\newblock \emph{Winning Ways for Your Mathematical Plays}.
\newblock Academic Press, 1982. (Second edition, A K Peters, 2001--2004.)

\bibitem{BW1994}
E.~Berlekamp and D.~Wolfe.
\newblock \emph{Mathematical Go: Chilling Gets the Last Point}.
\newblock A K Peters, 1994.

\bibitem{Conway1976}
J.~H. Conway.
\newblock \emph{On Numbers and Games}.
\newblock Academic Press, 1976. (Second edition, A K Peters, 2001.)

\bibitem{CrasmaruTromp2000}
M.~Cr\^{a}\c{s}maru and J.~Tromp.
\newblock Ladders are {PSPACE}-complete.
\newblock In \emph{Proceedings of the 2nd International Conference on Computers
  and Games (CG~2000)}, Lecture Notes in Computer Science, vol.~2063,
  pp.~241--249. Springer, 2001.

\bibitem{Robson1983}
J.~M. Robson.
\newblock The complexity of {G}o.
\newblock In \emph{Information Processing 83 (Proc. IFIP World Computer
  Congress)}, pp.~413--417. North-Holland, 1983.

\bibitem{Siegel2013}
A.~N. Siegel.
\newblock \emph{Combinatorial Game Theory}.
\newblock Graduate Studies in Mathematics, vol.~146. American Mathematical
  Society, 2013.

\bibitem{Spight2002}
W.~L. Spight.
\newblock Go thermography: the 4/21/98 Jiang--Rui endgame.
\newblock In R.~J. Nowakowski, editor, \emph{More Games of No Chance},
  MSRI Publications, vol.~42, pp.~89--105. Cambridge University Press, 2002.

\end{thebibliography}

\end{document}
